\documentclass[11pt,a4paper]{article}

\usepackage[utf8]{inputenc}
\usepackage[T1]{fontenc}
\usepackage[french]{babel}
\usepackage[a4paper,margin=2.2cm]{geometry}
\usepackage{amsmath,amssymb}
\usepackage{lmodern}
\usepackage{enumitem}
\usepackage{xcolor}
\usepackage{fancyhdr}
\usepackage{lastpage}
\usepackage{titlesec}
\usepackage{tcolorbox}
\usepackage{array}
\usepackage{tabularx}

\pagestyle{fancy}
\fancyhf{}
\lhead{\textbf{DM0 de mathématiques -- Niveau 4e}}
\rhead{\textbf{Nom : \hspace{4cm}}}
\cfoot{\thepage/\pageref{LastPage}}

\setlength{\parindent}{0pt}
\setlength{\parskip}{0.4em}

\definecolor{bleufonce}{RGB}{25,60,120}
\definecolor{bleuclair}{RGB}{235,242,255}

\titleformat{\section}
{\normalfont\Large\bfseries\color{bleufonce}}{}{0em}{}

\tcbset{
    colback=bleuclair,
    colframe=bleufonce,
    arc=2mm,
    boxrule=0.8pt
}

\begin{document}

\begin{center}
    {\LARGE \textbf{DM0 de mathématiques}}\\[0.4em]
    {\large Niveau 4e}\\[0.8em]
    \rule{0.9\textwidth}{0.8pt}
\end{center}

\begin{tabularx}{\textwidth}{|>{\centering\arraybackslash}p{0.22\textwidth}|>{\centering\arraybackslash}p{0.22\textwidth}|>{\centering\arraybackslash}X|}
\hline
\textbf{Durée} & \textbf{Barème} & \textbf{Matériel autorisé} \\
\hline
1 h 30 & 20 points & Calculatrice autorisée \\
\hline
\end{tabularx}

\vspace{1em}

\begin{tcolorbox}[title=Consignes générales]
\begin{itemize}[leftmargin=1.2cm]
    \item La rédaction compte dans tous les exercices.
    \item Chaque réponse doit être justifiée.
    \item Il faut écrire les calculs et faire des phrases de conclusion.
    \item Tout résultat non évident doit pouvoir être expliqué clairement à l'oral.
\end{itemize}
\end{tcolorbox}

\vspace{1em}

\begin{center}
\renewcommand{\arraystretch}{1.4}
\begin{tabular}{|c|c|c|c|c|}
\hline
\textbf{Exercice 1} & \textbf{Exercice 2} & \textbf{Exercice 3} & \textbf{Exercice 4} & \textbf{Total} \\
\hline
5 pts & 5 pts & 5 pts & 5 pts & 20 pts \\
\hline
\end{tabular}
\end{center}

\vspace{1em}

\section*{Exercice 1 -- Calcul littéral et raisonnement \hfill /5}

On considère l'expression :
\[
A = 3x + 5 - (2x - 1)
\]

\begin{enumerate}[label=\textbf{\arabic*.}, leftmargin=1.2cm]
    \item Réduire l'expression \(A\).
    \item Calculer \(A\) pour \(x=4\).
    \item Lina affirme que \(A = x+4\). A-t-elle raison ? Justifier.
    \item Déterminer la valeur de \(x\) pour laquelle \(A=10\).
\end{enumerate}

\begin{tcolorbox}[title=Rappel utile]
Réduire une expression signifie regrouper les termes de même nature.
\end{tcolorbox}

\textbf{Barème indicatif :} 1,5 pt ; 1 pt ; 1 pt ; 1,5 pt.

\section*{Exercice 2 -- Démonstration avec des triangles \hfill /5}

On considère un triangle \(ABC\) isocèle en \(A\), c'est-à-dire tel que :
\[
AB = AC.
\]
On note \(M\) le milieu du segment \([BC]\).

On veut démontrer que la droite \((AM)\) partage le triangle en deux triangles égaux.

\begin{tcolorbox}[title=Rappels utiles]
\begin{itemize}[leftmargin=1.2cm]
    \item Dire que \(M\) est le milieu de \([BC]\), c'est dire que \(BM = MC\).
    \item Deux triangles sont égaux s'ils ont leurs trois côtés de même longueur.
\end{itemize}
\end{tcolorbox}

\begin{enumerate}[label=\textbf{\arabic*.}, leftmargin=1.2cm]
    \item Expliquer pourquoi \(BM = MC\).
    \item Comparer les triangles \(ABM\) et \(ACM\).
    \item Montrer que les triangles \(ABM\) et \(ACM\) sont égaux.
    \item Écrire une phrase de conclusion expliquant ce que représente la droite \((AM)\) dans ce triangle.
\end{enumerate}

\textbf{Barème indicatif :} 1 pt ; 1 pt ; 2 pts ; 1 pt.

\section*{Exercice 3 -- Angles et schéma \hfill /5}

\textbf{Consigne :} \emph{À l'aide d'un schéma, démontrer que la somme des angles d'un triangle est égale à \(180^\circ\).}

On considère un triangle \(ABC\).

\begin{enumerate}[label=\textbf{\arabic*.}, leftmargin=1.2cm]
    \item Faire un schéma propre du triangle \(ABC\).
    \item Par le point \(A\), tracer une droite parallèle à la droite \((BC)\).
    \item En utilisant le schéma, expliquer pourquoi la somme des trois angles du triangle vaut \(180^\circ\).
\end{enumerate}

\begin{tcolorbox}[title=Attendus]
On attend :
\begin{itemize}[leftmargin=1.2cm]
    \item un schéma clair ;
    \item des angles nommés correctement ;
    \item une explication rédigée avec le mot \textbf{parallèle} ;
    \item une conclusion complète.
\end{itemize}
\end{tcolorbox}

\textbf{Barème indicatif :} 1 pt ; 1 pt ; 3 pts.

\section*{Exercice 4 -- Théorème de Pythagore et raisonnement \hfill /5}

On considère un triangle \(ABC\) rectangle en \(A\) tel que :
\[
AB = 6\ \text{cm} \qquad AC = 8\ \text{cm}
\]

\begin{enumerate}[label=\textbf{\arabic*.}, leftmargin=1.2cm]
    \item Calculer \(AB^2\) et \(AC^2\).
    \item En utilisant le théorème de Pythagore, calculer \(BC^2\).
    \item En déduire la longueur \(BC\).
    \item Sami affirme : \og Le triangle a pour côtés 6 cm, 8 cm et 10 cm \fg. A-t-il raison ? Justifier.
\end{enumerate}

\begin{tcolorbox}[title=Rappel utile]
Dans un triangle rectangle, le carré de l'hypoténuse est égal à la somme des carrés des deux autres côtés :
\[
BC^2 = AB^2 + AC^2
\]
si le triangle est rectangle en \(A\).
\end{tcolorbox}

\textbf{Barème indicatif :} 1 pt ; 1 pt ; 2 pts ; 1 pt.

\vfill

\begin{center}
\rule{0.8\textwidth}{0.6pt}\\[0.4em]
\textit{Fin du sujet -- La qualité du raisonnement est prise en compte.}
\end{center}

\end{document}